\documentclass{beamer}
\usetheme{AnnArbor}
\usepackage{apacite}
\setbeamertemplate{headline}{}
\title{FACTOR ANALYSIS}
\subtitle{LITERATURE REVIEW}
\date{MAY 2023}

\begin{document}
	\maketitle
	\begin{frame}
		In this chapter, we review the existing research and scholarship related to mental health and its effects on a person’s productivity as well as the interaction between an individual’s cognitive, affective and somatic symptoms of depression. A description of this consensus inevitably leads to a set of core issues that must be addressed if psychiatric scientists are to establish a more adequate explanation of mental health assessment and procedures. Several questions need to be answered with respect to; How an individual and environment contribute to a person’s mental health? How does mental health affect productivity at workplaces as well as learning institutions? How does mental health vary across age groups with the changing world of internet error post-COVID-19? What research methods are available for addressing these issues in research? The study begins with some key concepts from previous research perspectives that have contributed to recent advances in mental health understanding and then presents the implications of this knowledge for accurate screening of major depressive disorder.
	\end{frame}
\begin{frame}
	In a research done by \cite{martin2009meta}, investigated whether different types of health promotion intervention in the workplace reduce depression and anxiety symptoms. A systematic review and meta-analysis were undertaken on workplace health promotion published during the period 1997-2007.Studies were considered eligible for inclusion if they evaluated the impact of an intervention using a valid indicator or specific measure of depression or anxiety symptoms. The standardized mean difference was calculated for each of the following three types of outcome measures; depression, anxiety and composite mental health. Altogether,22 studies were found that met the inclusion criteria, with a total sample size of 3409 employees post intervention, and 17 of these studies were included in the meta-analysis representing 20 intervention- control comparisons. The pooled results indicated small but positive overall effects of the interventions with respect to symptoms of depression and anxiety, but no effect on composite mental health measures. The interventions that included a direct focus on mental health had a comparable effect on depression and anxiety symptoms as did the interventions with an indirect focus on risk factors. When the aim to reduce symptoms of depression and anxiety in employee populations, a broad range of health promotion interventions appear to be effective, although the effect is small.
\end{frame}
\begin{frame}
	 \cite{jain2013patient} examined the burden of depression on work productivity in the US workforce and used Patient Health Questionnaire for depression severity, the Health and Work Performance Questionnaire and Work Productivity and Activity Impairment (WPAI) Questionnaire for absenteeism and presenteeism in 1051 full-time employees for a period of three years. All levels of depression were associated with decreased work productivity and presenteeism was positively associated with severity of depression as well absenteeism which was significantly positively associated with severity of depression using the WPAI. The study showed that decreased overall productivity at all levels of depression and as severity increased, presenteeism and absenteeism worsened.
\end{frame}
\begin{frame}
	In another study \cite{serpa2014mindfulness}, investigated major problems among Veterans specifically anxiety, depression and suicidal ideation and used pre and post Mindfulness-based Stress Reduction questionnaires to collect data on 79 veterans at an urban Veterans Health Administration medical facility and also conducted a meditation analysis to determine whether changes in mindfulness were related to changes in the other outcomes for a period of 9 weekly sessions. This led to significant reductions in anxiety, depressions, and suicidal ideation after MBSR training, improved mental health functioning scores although pain intensity and physical health functionality did not indicate improvements. The study showed that completing an MBSR program can improve symptoms of anxiety, depression and reduced suicidal ideations on health of patients.
\end{frame}
\begin{frame}
	In the changing world, there is a need to assess workplace health culture using evidence-based measures. In a research done by \cite{laing2016anxiety}, indicated that previous studies using employee level culture of health measures observed that culture of health is associated with higher levels of job satisfaction and retention [15], work engagement [16], and healthy behaviors such as physical activity and healthy eating [20]. A 2016 on US government employees (n = 4703) found that workplace culture of health was negatively correlated with anxiety and depression [21]. Similar results were observed in employees in China, where higher ratings of workplace culture of health were associated with better psychological wellbeing [22]. 
	Domains of workplace culture of health such as leadership and coworkers support have been found to decrease employee health risk [8] and promote positive health behaviors [20]. However, there is only one U.S.-based study that has examined workplace health culture and emotional wellbeing from employees' perspective [21].  Other studies in this area have measured workplace culture of health from the employer-level, which may not capture a comprehensive perspective [18,19]. In addition, that study did not use a validated research tool to measure workplace culture of health, limiting its generalizability [21]. There is a need to assess workplace health culture using evidence-based measures. 
\end{frame}
\begin{frame}
	Another study was done by \cite{magtibay2017decreasing} to assess efficacy of blended learning to decrease stress and burnout among nurses through use of the Stress Management and Resiliency Training (SMART) program. It was assumed that Job-related stress in nurses leads to high rates of burnout, compromises patient care, and costs US healthcare organizations billions of dollars annually. Many mindfulness and resiliency programs are taught in a format that limits nurses’ attendance. Consistent with blended learning, participants chose the format that met their learning styles and goals; Web-based, independent reading, facilitated discussions. The end points of mindfulness, resilience, anxiety, stress, happiness, and burnout were measured at baseline, postintervention, and 3-month follow-up to examine within-group differences.Findings showed statistically significant, clinically meaningful decreases in anxiety, stress, and burnout and increases in resilience, happiness, and mindfulness. Results support blended learning using SMART as a strategy to increase access to resiliency training for nursing staff.
\end{frame}
\begin{frame}
	 In recent years, mental health has widely been recognized as the prevalence of mental health disorders continues to rise. In this regard there has been need to come up with screening techniques to ensure accurate and effective screening tools for mental health disorders. PHQ-9 is one the tools that have been employed over time to screen severity of depression. Patient Health Questionnaire (PHQ-9) is a brief, self-administered screening tool that is commonly used to diagnose and assess the severity of depression. The questionnaire consists of nine items that correspond to the nine DSM-5 criteria for major depressive disorder. Each item asks the patient to rate the frequency of a particular symptom over the past two weeks on a scale from 0 (not at all) to 3 (nearly every day). The scores are then summed to give a total score that ranges from 0 to 27, with higher scores indicating greater severity of depression. \cite{levis2019accuracy}, evaluates the accuracy of the Patient Health Questionnaire-9 (PHQ-9) for screening major depression in adults using individual participant data meta-analysis. The study, which involved the DEPRESSD Collaboration, analyzed data from 27 studies, with a total of 50,788 participants. The research found that a cut-off score of 10 or above can be used regardless of age, and the PHQ-9 is similarly sensitive but may be less specific for younger patients than for older patients.
\end{frame}
\begin{frame}
	The study also found that the PHQ-9 sensitivity compared with semi structured diagnostic interviews was greater than in previous conventional meta-analyses that combined reference standards and suggest that future research on the PHQ-9 should ideally be based on semi structured diagnostic interviews and should consider estimating probabilities of depression across the full spectrum of PHQ-9 screening scores, and should combine screening scores with individual characteristics to generate individualized probabilities of major depression. The research developed a web-based tool to estimate the expected number of positive screens and true and false screening outcomes based on the study's results. However, in primary care, approximately half of patients with positive screens would be false positives if this was used in practice.
\end{frame}
\begin{frame}
	\cite{leung2020measurement} aims to evaluate the factor structure, reliability, and construct validity of the Patient Health Questionnaire-9 (PHQ-9) in a community sample of Chinese adolescents. Confirmatory factor analysis was used to examine the one-factor structure of the PHQ-9, and multiple-group CFAs were performed to test for measurement invariance across gender and age subgroups. Internal consistency reliability was evaluated using Cronbach alpha, and construct validity was assessed using Pearson correlation coefficients between PHQ-9 scores and relevant measures. The results indicate that the PHQ-9 has a one-factor structure, good reliability, and construct validity in Chinese adolescents, and it is measurement invariant across gender and age subgroups. The study's limitations include the lack of test-retest reliability evaluation and generalizability to other cultural contexts. These findings provide important information on the psychometric properties of the PHQ-9 for assessing adolescent depression and screening, highlighting the need for further research on its predictive power and generalizability.
\end{frame}
\begin{frame}
In a more recent study, \cite{lister2023mental},investigated the  barriers and enablers to mental wellbeing and study success that students experienced in distance learning. Sixteen distance learning and five students were interviewed using narrative enquiry; students told their own stories and tutors told stories of students they had supported. A wide range of barriers and enablers were identified in the interviews, and these could be mapped to a number of areas in the university or study experience.As there were relationships between themes (i.e between assessment and study skills,or between study spaces and curriculum). The barriers and enablers are represented as a taxonomic wheel which classifies where barriers reside within distance learning,it maps them to diametrically corresponding enablers and indicates relationships between adjacent barriers and enablers.This taxonomy clearly demonstrated that  both barriers and enablers to student mental wellbeing in study reside throughout learning environments and study experiences as well as within the skillsets and capabilities of the students and supports the contention that holistic approaches are needed to support student mental wellbeing. A key finding of this study is that barriers and enablers were found in a majority of themes , implying that aspects of the study can be both barriers or enablers for mental wellbeing and study success, depending on who is experiencing them and how. The study indicated that therapeutic and individualistic approaches to mental health while important,are not sufficient to facilitate student mental wellbeing. The findings suggested that universities needed to tackle barriers where they reside, including exploring inclusive design practices, for curricula and assessment, explicitly teaching study skills,designing systems, and processes that do not cause undue stress and designing learning spaces to be communities.
\end{frame}
\begin{frame}
	By adopting a holistic approach to mental health in this way, distance learning can begin to be truly conducive to student wellbeing.  The existence of the gap to fully explore the relationship between an individual's cognitive, affective and somatic symptoms of depression, decision to advance the research has been put into consideration. Different methods of CFA have been used to analyze and investigate if there is existence of any kind. In our research, multiple group confirmatory factor analysis is the key interest in observing this relationship.
\end{frame}
\begin{frame}
	\bibliographystyle{apacite}
	\bibliography{Refferences}
\end{frame}
\end{document}